\section{Analysis of Model-\texorpdfstring{$A$}{A} (\texorpdfstring{$\gamma_1=\gamma_2=0$}{gamma1=gamma2=0}, \texorpdfstring{$\theta_1=\theta_2=0$}{theta1=theta2=0})}\label{sec:model_a}

Model-$A$ is the baseline of this collection: all jockeying and abandonment
parameters are zero, leaving an ordinary two-class $M/M/1$ queue with
non-preemptive priority. We analyse it in both state spaces $S$ and
$\widetilde{S}$: on the former we give two complete proofs, on the latter
a not very accurate approximate result, for which an evaluation is provided in Approximation~\ref{apx:A:PPGF_Stilde}.

The balance equations are obtained by plugging $\gamma_1=\gamma_2=\theta_1=\theta_2=0$ into Equations~\eqref{eq:gen:balance_interior},
\eqref{eq:gen:balance_x_boundary}, \eqref{eq:gen:balance_y_boundary},
\eqref{eq:gen:balance_zero}, and~\eqref{eq:gen:balance_idle} for state space $S$, and
Equations~\eqref{eq:gen:tilde_balance_interior},
\eqref{eq:gen:tilde_balance_diagonal}, \eqref{eq:gen:tilde_balance_y_boundary},
\eqref{eq:gen:tilde_balance_zero}, and~\eqref{eq:gen:tilde_balance_idle} for state
space $\widetilde{S}$. The corresponding starting points are
Lemma~\ref{lem:gen:fundamental_equation} for $S$ and
Lemma~\ref{lem:gen:PPGF_dynamics} for $\widetilde{S}$.

\textbf{Stability.} Since there is no abandonment, the system is positive recurrent if
and only if the total load satisfies
\begin{equation}
    \rho = \frac{\lambda_1+\lambda_2}{\mu} < 1.
    \label{eq:A:stability}
\end{equation}

Figure~\ref{fig:model-a-queue} shows the queueing diagram of Model-$A$.
\begin{figure}[H]
    \centering
    \resizebox{0.62\textwidth}{!}{%
        % Model A: γ₁=γ₂=θ₁=θ₂=0
% Two queues, one server. Arrivals and service only — no jockeying, no abandonment.
\begin{tikzpicture}[>=Latex, thick]

    % ── Queue 1 (top, Class-1, priority) ──────────────────────────────────────
    \draw (0, 1.9) -- (2, 1.9) -- (2, 1.1) -- (0, 1.1);
    \foreach \x in {0.2, 0.4, ..., 1.8} { \draw[thin] (\x, 1.8) -- (\x, 1.2); }
    \draw[->] (-1.5, 1.5) -- (-0.1, 1.5) node[midway, above] {$\lambda_1$};

    % ── Queue 2 (bottom, Class-2) ──────────────────────────────────────────────
    \draw (0, -1.1) -- (2, -1.1) -- (2, -1.9) -- (0, -1.9);
    \foreach \x in {0.2, 0.4, ..., 1.8} { \draw[thin] (\x, -1.2) -- (\x, -1.8); }
    \draw[->] (-1.5, -1.5) -- (-0.1, -1.5) node[midway, above] {$\lambda_2$};

    % ── Server ────────────────────────────────────────────────────────────────
    \node[draw, circle, minimum size=1.2cm] (server) at (5, 0) {$\mu$};
    \draw[->] (2.2,  1.5) -- (server);
    \draw[->] (2.2, -1.5) -- (server);
    \draw[->] (server) -- (7, 0);

\end{tikzpicture}
%
    }
    \caption{Queue layout for Model-$A$. Only Poisson arrivals and exponential service
    are present; all active mechanisms rates are zero.}
    \label{fig:model-a-queue}
\end{figure}

\subsection{Analysis of Model-\texorpdfstring{$A$}{A}: state space \texorpdfstring{$S$}{S}}
We begin from Equation~\eqref{eq:gen:fundamental_equation} with
$\gamma_1=\gamma_2=\theta_1=\theta_2=0$, in which the marginal $P_y(y)=P(0,y)$ is
still to be determined:
\begin{align}
    \begin{split}
        \left[ (\lambda_1+\lambda_2+ \mu)xy - \mu y - \lambda_1 x^2 y - \lambda_2 x y^2 \right] P(x,y)
         &= \mu(x-y) P_y(y)
         - \mu x(1-y) \pi(0,0).
    \end{split} \label{eq:A:fundamental}
\end{align}

The following theorem summarises the main result of this subsection; two standalone
proofs are given, one analytical and one probabilistic.
\begin{theorem}\label{thm:A:PGF}
    Consider a queueing system with two priority classes with Poisson arrival rates
    $\lambda_1$ and $\lambda_2$, where class-$1$ customers have non-preemptive priority
    over those of class~$2$, and exponentially distributed service times with rate $\mu$
    independently of the class. Under the stability condition $\rho<1$~\eqref{eq:A:stability}, the
    joint PGF $P(x,y) = \mathbb{E}[x^{N_1}y^{N_2}]$ is
    \begin{align}
        P(x,y) &= \left[ (1+\rho_1 +\rho_2)xy - y -\rho_1x^2y - \rho_2xy^2 \right]^{-1}
                  \cdot y(1-y) \left[ \frac{x - x^*(y)}{x^*(y)-y} \right] \cdot \rho (1-\rho),
        \label{eq:A:PGF}
    \end{align}
    where
    \begin{equation}
        x^*(y) = \frac{\left[\mu + \lambda_1 + \lambda_2 (1-y)\right]
                       - \sqrt{\left[\mu + \lambda_1 + \lambda_2 (1-y)\right]^2 - 4 \lambda_1 \mu}}
                      {2 \lambda_1}.
        \label{eq:A:xstar_def}
    \end{equation}
\end{theorem}

\begin{corollary}\label{cor:A:pi0}
    Under the stability condition $\rho < 1$~\eqref{eq:A:stability}, the empty-state probabilities
    for Model-$A$ are
    \begin{equation}
        \pi_0 = 1-\rho, \qquad \pi(0,0) = \rho(1-\rho). \label{eq:A:pi0_pi00}
    \end{equation}
    In particular, $P(1,1) = 1-\pi_0 = \rho$ and the diagonal generating function
    satisfies $P(z,z)=\pi(0,0)/(1-\rho z)$.
\end{corollary}

\begin{proof}[Proof of Corollary~\ref{cor:A:pi0}]
Setting $x=y=z$ in the fundamental equation~\eqref{eq:A:fundamental}, the boundary
term $\mu(x-y)P_y(y)$ vanishes:
    \begin{align*}
        \left[ (\lambda_1+\lambda_2+ \mu)z^2 - \mu z - \lambda_1 z^3 - \lambda_2 z^3 \right] P(z,z)
        &= - \mu z(1-z) \pi(0,0).
    \end{align*}
    Dividing both sides by $z$ and factoring the bracket:
    \begin{align*}
        -(1-z)\left[ \mu - (\lambda_1 + \lambda_2)z \right] P(z,z)
        &= - \mu (1-z) \pi(0,0).
    \end{align*}
    Cancelling $(1-z)\neq 0$ for $z\in[0,1)$ and rearranging:
    \begin{equation*}
        P(z,z) = \frac{\mu\,\pi(0,0)}{\mu - \lambda z} = \frac{\pi(0,0)}{1-\rho z},
    \end{equation*}
    where $\lambda = \lambda_1+\lambda_2$ and the second equality follows by dividing
    numerator and denominator by $\mu$.

    Setting $z=1$ and using the normalisation $P(1,1) = 1-\pi_0$:
    \begin{equation}
        1-\pi_0 = \frac{\pi(0,0)}{1-\rho}. \label{eq:A:norm_step}
    \end{equation}
    The global balance equation for the idle state gives $\lambda\pi_0 = \mu\pi(0,0)$,
    i.e.\ $\pi(0,0) = \rho\pi_0$. Substituting into~\eqref{eq:A:norm_step}:
    \begin{equation*}
        1-\pi_0 = \frac{\rho\pi_0}{1-\rho}
        \implies (1-\pi_0)(1-\rho) = \rho\pi_0
        \implies \pi_0 = 1-\rho,
    \end{equation*}
    and therefore $\pi(0,0) = \rho(1-\rho)$.
\end{proof}

\subsubsection{Analytical approach for Model-\texorpdfstring{$A$}{A} in state space \texorpdfstring{$S$}{S}}\label{sssec:A:analytical}

\begin{proof}[Proof of Theorem~\ref{thm:A:PGF}, analytical approach.]
    We begin from the fundamental equation for Model-$A$~\eqref{eq:A:fundamental}:
    \begin{equation*}
        \left[ (\lambda_1+\lambda_2+ \mu)xy - \mu y - \lambda_1 x^2y - \lambda_2 xy^2 \right] P(x,y)
        = \mu(x-y) P_y(y)  - \mu x(1-y) \pi(0,0).
    \end{equation*}
    We apply the \emph{Kernel Method}, used in ~\cite{cohen1956delay}: fix $y\in(0,1]$ and find $x^*(y)$ such that the LHS of the expression above is zero. Since $P(x,y)$ is a PGF, it is bounded. At any such root the LHS vanishes, so the RHS must also vanish:
    \[
        \mu\bigl[(x^*(y)-y)\,P_y(y) - x^*(y)(1-y)\,\pi(0,0)\bigr] = 0.
    \]
    To find the roots of the coefficient of $P(x,y)$, divide by $y\neq 0$ and set
    \begin{align*}
        \begin{split}
             0 &=  (\lambda_1 + \lambda_2 + \mu)x - \mu  -\lambda_1x^2 - \lambda_2xy
        \end{split} \\
        \begin{split}
          0 &=  \lambda_1 x^2 - \left[ \mu + \lambda_1 + \lambda_2 (1-y) \right] x + \mu,
        \end{split}
    \end{align*}
    whose roots are
    \begin{equation*}
        x^{\pm}(y) = \frac{
            \left[\mu + \lambda_1 + \lambda_2 (1-y) \right] \pm \sqrt{
                \left[\mu + \lambda_1 + \lambda_2 (1-y) \right]^2 - 4  \mu \lambda_1
            }
        }{
            2\lambda_1
        }.
    \end{equation*}
    
    We label the root with the negative sign $x^*(y)$ and the root with the positive
    sign $x^+(y)$. At $y=1$:
    \begin{equation*}
        x^*(1) = 1 \quad \text{and} \quad x^+(1) = \mu/\lambda_1 > 1.
    \end{equation*}
    
    We require our root to lie inside the closed unit disk for $y\in[0,1)$. Computing
    the derivative of $x^*(y)$:
    \begin{equation}
        \frac{d}{dy}x^*(y) = \frac{\lambda_2}{2\lambda_1}\!\left[-1+\frac{\mu+\lambda_1+\lambda_2(1-y)}{\sqrt{(\mu+\lambda_1+\lambda_2(1-y))^2-4\lambda_1\mu}}\right]>0,
        \label{eq:A:xstar_prime}
    \end{equation}
    since $(\mu+\lambda_1+\lambda_2(1-y))^2>(\mu+\lambda_1+\lambda_2(1-y))^2-4\lambda_1\mu$
    whenever $\lambda_1\mu>0$, by squaring numerator and denominator in the inner fraction. So $x^*(y)$ is strictly increasing in $[0,1]$, and since $x^*(1)=1$ we have $x^*(y)<1$ for $y\in[0,1)$.
    We make use of Vieta's formulas to discard the other root for the kernel polynomial
    $\lambda_1 x^2-[\mu+\lambda_1+\lambda_2(1-y)]x+\mu=0$: the product of the two
    roots satisfies $x^*(y)\,x^+(y)= \mu/\lambda_1>1$, so
    $x^+(y)=\mu/(\lambda_1 x^*(y))>1$ throughout $[0,1)$. So the only root
    strictly inside the unit disk for $y\in[0,1)$ is $x^*(y)$:
    \begin{equation}
        x^*(y) = \frac{
            (\mu +\lambda_1 + \lambda_2 (1-y))
            - \sqrt{(\mu +\lambda_1 + \lambda_2 (1-y))^2 - 4 \lambda_1 \mu}
        }{
            2\lambda_1
        }.
        \label{eq:A:xstar}
    \end{equation}

    Returning to the fundamental equation~\eqref{eq:A:fundamental} at $x=x^*(y)$,
    the RHS also equals zero, giving
    \begin{equation}
      \mu\cdot\left[(x^*(y)-y)P_y(y) - x^*(y)(1-y)\pi(0,0) \right]=0 \quad \implies \quad  P_y(y) = \frac{x^*(y)(1-y)}{x^*(y) - y} \pi(0,0). \label{eq:A:Py}
    \end{equation}

    By Corollary~\ref{cor:A:pi0}, $\pi(0,0)=\rho(1-\rho)$. Substituting this and
    Equation~\eqref{eq:A:Py} into the fundamental equation~\eqref{eq:A:fundamental}:
    \begin{align*}
        \begin{split}
            \left[ (\lambda_1+\lambda_2+ \mu)xy - \mu y - \lambda_1 x^2y - \lambda_2 xy^2 \right] P(x,y)
            &= \mu(x-y) \frac{x^*(y)(1-y)}{x^*(y) - y} \rho(1-\rho)  - \mu x(1-y) \rho(1-\rho)
        \end{split} \\
        \begin{split}
            \left[ (\lambda_1+\lambda_2+\mu)xy -\mu y-\lambda_1 x^2y - \lambda_2 xy^2 \right] P(x,y)
            &= \mu (1-y) \left[ \frac{x^*(y)(x-y)}{x^*(y)-y} - x \right] \rho (1-\rho)
        \end{split} \\
        \begin{split}
            \left[ (\lambda_1+\lambda_2+\mu)xy -\mu y-\lambda_1 x^2y - \lambda_2 xy^2 \right] P(x,y)
            &= \mu  (1-y) \left[ \frac{y(x - x^*(y))}{x^*(y) - y} \right] \rho (1-\rho)
        \end{split} \\
        \begin{split}
            \left[ (1 + \rho_1 + \rho_2)xy - y-\rho_1 x^2y - \rho_2 xy^2 \right] P(x,y)
            &= y (1-y) \left[ \frac{x - x^*(y)}{x^*(y) - y} \right] \rho (1-\rho).
        \end{split}
    \end{align*}

    Multiplying both sides by $\left[(1 + \rho_1 + \rho_2)xy - y-\rho_1 x^2y - \rho_2 xy^2\right]^{-1}$
    recovers $P(x,y)$ in Theorem~\ref{thm:A:PGF} and concludes the proof.
\end{proof}


\subsubsection{Probabilistic approach for Model-\texorpdfstring{$A$}{A} in state space \texorpdfstring{$S$}{S}}\label{sssec:A:probabilistic}
\begin{proof}[Proof of Theorem~\ref{thm:A:PGF}, probabilistic approach.]
    This proof mirrors the analytical one, but derives the class-$2$ marginal PGF $P_y(y)$ by means of probabilistic arguments. We again begin from the fundamental equation for Model-$A$ (Equation~\eqref{eq:A:fundamental}):
    \begin{equation}
        \left[ (\lambda_1 + \lambda_2 + \mu) xy - \mu y - \lambda_1 x^2 y - \lambda_2 x y^2 \right] P(x,y) =
        \mu \big[ (x-y) P_y(y) - x(1-y) \pi(0,0) \big].
        \label{eq:A:prob_starting}
    \end{equation}

    \paragraph{Class-$2$ conditional distribution via $M/G/1$.}
    Since class~$1$ has non-preemptive priority, class-$2$ customers are not taken into service
    whenever the priority queue is not empty ($N_1>0$); their \emph{effective service time} $B$ is therefore the duration of a complete busy period of the single-class ordinary $M/M/1$ sub-queue with arrival rate $\lambda_1$ and service rate $\mu$. From class-$2$'s viewpoint the system is thus an $M/G/1$ queue with arrival rate $\lambda_2$ and generally distributed service times $B$ and mean $\mathbb{E}[B]$. The LST and mean of $B$ are presented
    in~\S\ref{ssec:mm1} (Equations~\eqref{eq:mm1:bp_lst}--\eqref{eq:mm1:bp_mean}), now
    with $\lambda_1$ instead of the generic $\lambda$:
    \begin{align*}
        \widetilde{B}(s) &= \frac{(\mu+\lambda_1+s) - \sqrt{(\mu + \lambda_1 + s)^2-4\mu\lambda_1}}{2\lambda_1}, \\
        \mathbb{E}[B] &= \frac{1}{\mu(1-\rho_1)}. 
    \end{align*}
    Applying the Pollaczek--Khinchine (PK) formula (\S\ref{ssec:pk},
    Equation~\eqref{eq:pk:pgf}) to this $M/G/1$ system, with
    $\widetilde{B}(\lambda_2(1-y))=x^*(y)$ (cf.~\eqref{eq:A:xstar}) and
    $\lambda_2\mathbb{E}[B]=\rho_2/(1-\rho_1)$:
    \begin{align*}
        \begin{split}
            L_2(y) = \frac{(1 - \lambda_2 \mathbb{E}[B]) (1-y) \widetilde{B}(\lambda_2(1-y))}{\widetilde{B}(\lambda_2(1-y)) - y}
                    = \frac{1-\rho}{1-\rho_1} \cdot \frac{x^*(y)(1-y)}{x^*(y) - y}.
        \end{split}
    \end{align*}
    By the PASTA property, the time-average distribution of $N_2$ conditional on
    $(\text{busy},N_1=0)$ coincides with the $M/G/1$ steady-state distribution, so
    $\mathbb{E}[y^{N_2}\mid\text{busy},N_1=0]=L_2(y)$.

    By the law of total expectation, $P_y(y)=\mathbb{P}(\text{busy},N_1=0)\cdot\mathbb{E}[y^{N_2}\mid\text{busy},N_1=0]$. So
    \begin{equation*}
        P_y(y)=\mathbb{P}(\text{busy},N_1=0)\cdot\mathbb{E}[y^{N_2}\mid\text{busy},N_1=0] = \mathbb{P}(\text{busy},N_1=0)\cdot L_2(y).
    \end{equation*}
    Under non-preemptive priority, the priority queue operates autonomously (memorylessness of residual service time), so $\mathbb{P}(N_1=0 \mid \text{busy})=1-\rho_1$ by the geometric law of the ordinary $M/M/1(\lambda_1,\mu)$. Combined with $\mathbb{P}(\text{busy})=\rho$, we have
    \begin{equation*}
        \mathbb{P}(\text{busy},N_1=0)=\rho(1-\rho_1);
    \end{equation*}
    and therefore
    \begin{equation}
        P_y(y) = \rho(1-\rho_1)\cdot L_2(y) = \rho(1-\rho_1)\cdot\frac{1-\rho}{1-\rho_1} \cdot \frac{x^*(y)(1-y)}{x^*(y) - y} = \frac{x^*(y)(1-y)}{x^*(y) - y} \cdot (1-\rho) \rho
        \label{eq:A:Py_prob}
    \end{equation}

    Substituting~\eqref{eq:A:Py_prob} and $\pi(0,0)=\rho(1-\rho)$ (Corollary~\ref{cor:A:pi0})
    into~\eqref{eq:A:prob_starting}:
    \begin{align*}
        \begin{split}
            \left[ (\lambda_1+\lambda_2+\mu)xy - \mu y - \lambda_1 x^2y - \lambda_2 xy^2 \right] P(x,y)
            &= \mu(x-y) \frac{x^*(y)(1-y)}{x^*(y)-y} \rho(1-\rho) - \mu x(1-y) \rho(1-\rho)
        \end{split} \\
        \begin{split}
            \left[ (\lambda_1+\lambda_2+\mu)xy - \mu y - \lambda_1 x^2y - \lambda_2 xy^2 \right] P(x,y)&= \mu(1-y) \left[ \frac{x^*(y)(x-y)}{x^*(y)-y} - x \right] \rho(1-\rho)
        \end{split} \\
        \begin{split}
            \left[ (\lambda_1+\lambda_2+\mu)xy - \mu y - \lambda_1 x^2y - \lambda_2 xy^2 \right] P(x,y)&= \mu(1-y) \left[ \frac{y(x-x^*(y))}{x^*(y)-y} \right] \rho(1-\rho)
        \end{split} \\
        \begin{split}
            \left[ (1+\rho_1+\rho_2)xy - y - \rho_1 x^2y - \rho_2 xy^2 \right] P(x,y)
            &= y(1-y) \left[ \frac{x-x^*(y)}{x^*(y)-y} \right] \rho(1-\rho).
        \end{split}
    \end{align*}
    Multiplying both sides by $\left[ (1+\rho_1+\rho_2)xy - y - \rho_1 x^2y - \rho_2 xy^2 \right]^{-1}$
    to isolate $P(x,y)$ recovers Theorem~\ref{thm:A:PGF} and concludes the proof.
\end{proof}


\subsubsection{Partial-PGF recurrence approach for Model-\texorpdfstring{$A$}{A} in state space \texorpdfstring{$S$}{S}}


This proof follows Cohen~\cite{cohen1956delay}'s method. We keep the class-$1$
index $n_1$ as a discrete parameter and generate only with respect
to $N_2$, via the PPGF $\widetilde{P}(n_1,y)=\sum_{n_2\geq0}\pi(n_1,n_2)\,y^{n_2}$
from~\eqref{eq:gen:ppgf_n1}. Multiplying the balance equations by $y^{n_2}$ and
summing over $n_2$ turns them into a second-order linear
recurrence in $n_1$, whose characteristic roots fix the geometric structure of the
solution in the priority-class index.

\begin{proof}
    We start from the Model-$A$ balance equations, obtained by setting $\gamma_1=\gamma_2=\theta_1=\theta_2=0$ in the interior equation~\eqref{eq:gen:balance_interior} and the $x$-boundary equation~\eqref{eq:gen:balance_x_boundary}:
    \begin{align}
        (\lambda_1+\lambda_2+\mu)\,\pi(n_1,n_2)
            &= \mu\,\pi(n_1+1,n_2) + \lambda_1\,\pi(n_1-1,n_2) + \lambda_2\,\pi(n_1,n_2-1),
            && n_1\geq1,\ n_2\geq1, \nonumber\\
        (\lambda_1+\lambda_2+\mu)\,\pi(n_1,0)
            &= \mu\,\pi(n_1+1,0) + \lambda_1\,\pi(n_1-1,0),
            && n_1\geq1,\ n_2=0. \label{eq:A:ppgf_xbound}
    \end{align}
    Now we take transforms by summing all over their states, but only with respect to $n_2$, and we substitute our PPGF $\widetilde{P}(n_1,y)=\sum_{n_2=0}^{\infty} \pi(n_1,y)y^{n_2}$:
    \begin{align}
        \begin{split}
            & (\lambda_1+\lambda_2+\mu)\sum_{n_2=1}^{\infty}\pi(n_1,n_2)y^{n_2} \\
            & \quad = \quad \mu\sum_{n_2=1}^{\infty}\pi(n_1+1,n_2)y^{n_2} + \lambda_1\sum_{n_2=1}^{\infty}\pi(n_1-1,n_2)y^{n_2} + \lambda_2\sum_{n_2=1}^{\infty}\pi(n_1,n_2-1)y^{n_2}
        \end{split} \nonumber \\
        \begin{split}
            & LHS_1 \quad = \quad RHS_1 + RHS_2 + RHS_3
        \end{split} \nonumber \\
        \begin{split}
            & LHS_1 \\
            & \quad = \quad (\lambda_1+\lambda_2+\mu)\sum_{n_2=1}^{\infty}\pi(n_1,n_2)y^{n_2} \\
            & \quad = \quad (\lambda_1+\lambda_2+\mu) \left[ \sum_{n_2=0}^{\infty}\pi(n_1,n_2)y^{n_2} - \pi(n_1,0)y^0 \right] \\
            & \quad = \quad (\lambda_1+\lambda_2+\mu) \left[ \widetilde{P}(n_1,y) - \pi(n_1,0) \right]
        \end{split} \nonumber \\
        \begin{split}
            & RHS_1 \\
            & \quad = \quad \mu\sum_{n_2=1}^{\infty}\pi(n_1+1,n_2)y^{n_2} \\
            & \quad = \quad \mu \left[ \sum_{n_2=0}^{\infty}\pi(n_1+1,n_2)y^{n_2} - \pi(n_1+1,0)y^0 \right] \\
            & \quad = \quad \mu \left[ \widetilde{P}(n_1+1,y) - \pi(n_1+1,0) \right]
        \end{split} \nonumber \\
        \begin{split}
            & RHS_2 \\
            & \quad = \quad \lambda_1\sum_{n_2=1}^{\infty}\pi(n_1-1,n_2)y^{n_2} \\
            & \quad = \quad \lambda_1 \left[ \sum_{n_2=0}^{\infty}\pi(n_1-1,n_2)y^{n_2} - \pi(n_1-1,0)y^0 \right] \\
            & \quad = \quad \lambda_1 \left[ \widetilde{P}(n_1-1,y) - \pi(n_1-1,0) \right]
        \end{split} \nonumber \\
        \begin{split}
            & RHS_3 \\
            & \quad = \quad \lambda_2\sum_{n_2=1}^{\infty}\pi(n_1,n_2-1)y^{n_2} \\
            & \quad = \quad \lambda_2 y \sum_{n_2=1}^{\infty}\pi(n_1,n_2-1)y^{n_2-1} \\
            & \quad = \quad \lambda_2 y \sum_{m=0}^{\infty}\pi(n_1,m)y^{m} \\
            & \quad = \quad \lambda_2 y \, \widetilde{P}(n_1,y)
        \end{split} \nonumber \\
        \begin{split}
            & (\lambda_1+\lambda_2+\mu)\bigl[\widetilde{P}(n_1,y)-\pi(n_1,0)\bigr] \\
            & \quad = \quad \mu\bigl[\widetilde{P}(n_1+1,y)-\pi(n_1+1,0)\bigr] \\
            & \quad \quad + \; \lambda_1\bigl[\widetilde{P}(n_1-1,y)-\pi(n_1-1,0)\bigr] \\
            & \quad \quad + \; \lambda_2 y\,\widetilde{P}(n_1,y)
        \end{split} \nonumber \\
        \begin{split}
            & \bigl[\mu+\lambda_1+\lambda_2(1-y)\bigr]\widetilde{P}(n_1,y) -\mu\widetilde{P}(n_1+1,y)-\lambda_1\widetilde{P}(n_1-1,y) \\
            & \quad = \quad (\lambda_1+\lambda_2+\mu)\pi(n_1,0) -\mu\pi(n_1+1,0)-\lambda_1\pi(n_1-1,0)
        \end{split}
        \label{eq:A:ppgf_final_boundary}
    \end{align}
    By the $x$-boundary balance equation~\eqref{eq:A:ppgf_xbound}, the right-hand side of~\eqref{eq:A:ppgf_final_boundary} vanishes completely, giving the homogeneous recurrence
    \begin{equation}
        \left[ \mu + \lambda_1 + \lambda_2 (1-y) \right] \widetilde{P}(n_1,y)
        = \mu \widetilde{P}(n_1+1,y)
        + \lambda_1 \widetilde{P}(n_1-1,y),
        \qquad n_1\geq1.
        \label{eq:A:ppgf_recurrence}
    \end{equation}
    
    \paragraph{Characteristic roots and mode selection.}
    Equation~\eqref{eq:A:ppgf_recurrence} is a \emph{second-order linear recurrence} in $n_1$ with constant coefficients, so its general solution is a linear combination of two geometric sequences $[\widetilde{x}^{\pm}(y)]^{n_1}$, where $\widetilde{x}^{\pm}(y)$ solve the characteristic equation
    \begin{equation}
        \mu\,\widetilde{x}^2 - \left[ \mu + \lambda_1 + \lambda_2(1-y) \right]\widetilde{x} + \lambda_1 = 0,
        \label{eq:A:ppgf_char}
    \end{equation}
    that is,
    \begin{equation*}
        \widetilde{x}^{\pm}(y)
        = \frac{\bigl(\mu + \lambda_1 + \lambda_2(1-y)\bigr)
                \pm \sqrt{\bigl(\mu + \lambda_1 + \lambda_2(1-y)\bigr)^2 - 4\lambda_1\mu}}{2\mu}.
    \end{equation*}
    The \emph{characteristic polynomial} (Equation \eqref{eq:A:ppgf_char}) is the reciprocal to that of the bivariate PGF approaches (i.e. the reciprocal of $\lambda_1 x^2-[\mu+\lambda_1+\lambda_2(1-y)]x+\mu$), thus its roots are the reciprocals of that kernel's roots: the positive root is the reciprocal of the negative one ($\widetilde{x}^{+}(y)=1/x^*(y)$), and conversely $\widetilde{x}^{-}(y)=1/x^{+}(y)$. By Vieta's formulas, $\widetilde{x}^{-}(y)\,\widetilde{x}^{+}(y)=\lambda_1/\mu=\rho_1<1$. Recall from \S\ref{sssec:A:analytical} that $x^*(y)<1<x^{+}(y)$ for $y\in[0,1)$.
    We select the negative-sign root because the sequence $\{\widetilde{P}(n_1,y)\}_{n_1\geq0}$ must be summable:
    \begin{equation*}
        \sum_{n_1=0}^{\infty}\widetilde{P}(n_1,y) = P(1,y) \le P(1,1) = 1-\pi_0 = \rho < \infty.
    \end{equation*}
    At $y=1$, we have
    \begin{equation*}
        \widetilde{x}^{-}(1)=\rho_1 \quad \text{and} \quad \widetilde{x}^{+}(1)=1,
    \end{equation*}
    the latter making the geometric sequence $[\widetilde{x}^{+}(y)]^{n_1}$ non-summable. We therefore define $\widetilde{x}^*(y):=\widetilde{x}^{-}(y)$,
    \begin{equation}
        \widetilde{x}^*(y):=\widetilde{x}^{-}(y) = \frac{\bigl(\mu + \lambda_1 + \lambda_2(1-y)\bigr)
                - \sqrt{\bigl(\mu + \lambda_1 + \lambda_2(1-y)\bigr)^2 - 4\lambda_1\mu}}{2\mu}
        = \rho_1\,x^*(y).
        \label{eq:A:xtildestar}
    \end{equation}
    The last identity holds because $\widetilde{x}^*$ and $x^*$ share the same numerator and differ only in the denominator ($2\mu$ versus $2\lambda_1$). Next,
    \begin{equation}
        \widetilde{P}(n_1, y) = \widetilde{P}(0,y)\,[\widetilde{x}^*(y)]^{n_1},
        \qquad n_1\geq0,
        \label{eq:A:ppgf_geometric}
    \end{equation}
    where the relation also holds at $n_1=0$ trivially.
    Recall that at $y=1$, we have $\widetilde{x}^*(1)=\lambda_1/\mu=\rho_1$. Thus, $\widetilde{P}(n_1,1)=\widetilde{P}(0,1)\,\rho_1^{n_1}$ and $\widetilde{P}(0,1)=P_y(1)=\rho(1-\rho_1)$. The constant $\widetilde{P}(0,y)$ is resolved by its definition
    \begin{equation*}
        \widetilde{P}(0,y) = \sum_{n_2=0}^{\infty}\pi(0,n_2)\,y^{n_2} = P(0,y) = P_y(y),
    \end{equation*}
    as determined in the bivariate PGF approaches.
    Summing the geometric law (Equation~\eqref{eq:A:ppgf_geometric}) over $n_1 \geq 0$ for $x \in (0,1]$ recovers
    \begin{equation}
        P(x,y) = \sum_{n_1=0}^{\infty}\widetilde{P}(n_1,y)\,x^{n_1}
               = P_y(y)\sum_{n_1=0}^{\infty}\bigl[\widetilde{x}^*(y)\,x\bigr]^{n_1}
               = \frac{P_y(y)}{1-\widetilde{x}^*(y)\,x}.
        \label{eq:A:ppgf_reconstruct}
    \end{equation}
    To bring this to the form of Theorem~\ref{thm:A:PGF}, we define the kernel of the fundamental equation, divided by $\mu$:
    \begin{equation}
        K(x,y) := (1+\rho_1+\rho_2)xy - y - \rho_1 x^2 y - \rho_2 x y^2
                = -y\rho_1\bigl(x-x^*(y)\bigr)\bigl(x-x^{+}(y)\bigr),
        \label{eq:A:kernel_factored}
    \end{equation}
    where $x^*(y)$ and $x^{+}(y)$ are the roots found in \S\ref{sssec:A:analytical}. By Vieta's product formula, $x^*(y)\,x^{+}(y)=\mu/\lambda_1=1/\rho_1$. This, together with $\widetilde{x}^*(y)$ (Equation \eqref{eq:A:xtildestar}), gives $\widetilde{x}^*(y)=\rho_1 x^*(y)=1/x^{+}(y)$.
    Using the factorisation of $K(x,y)$~\eqref{eq:A:kernel_factored}, the denominator of the reconstructed PGF~\eqref{eq:A:ppgf_reconstruct} becomes
    \begin{equation*}
        1-\widetilde{x}^*(y)\,x = 1 - \frac{x}{x^{+}(y)} = \frac{x^{+}(y)-x}{x^{+}(y)}
        = \frac{K(x,y)}{y\rho_1\,x^{+}(y)\bigl(x-x^*(y)\bigr)}.
    \end{equation*}
    Substituting this into our reconstructed $P(x,y)$ (Equation \eqref{eq:A:ppgf_reconstruct}), using 
    $P_y(y)=\frac{x^*(y)(1-y)}{x^*(y)-y}\rho(1-\rho)$ and also noting that $\rho_1\,x^*(y)\,x^{+}(y)=1$, we obtain
    \begin{align*}
        P(x,y)
        &= \frac{x^*(y)(1-y)}{x^*(y)-y}\,\rho(1-\rho)
           \cdot \frac{y\rho_1\,x^{+}(y)\bigl(x-x^*(y)\bigr)}{K(x,y)} \\
        &= \underbrace{\rho_1\,x^*(y)\,x^{+}(y)}_{=\,1}\;
           \frac{1}{K(x,y)}\, y(1-y)\,\frac{x-x^*(y)}{x^*(y)-y}\,\rho(1-\rho) \\
        &= \bigl[(1+\rho_1+\rho_2)xy - y - \rho_1 x^2 y - \rho_2 x y^2\bigr]^{-1}
           \cdot y(1-y)\,\frac{x-x^*(y)}{x^*(y)-y}\,\rho(1-\rho),
    \end{align*}
    which recovers Equation~\eqref{eq:A:PGF} from Theorem~\ref{thm:A:PGF} and concludes the proof.
\end{proof}


\subsection{Analysis of Model-\texorpdfstring{$A$}{A} on \texorpdfstring{$\widetilde{S}$}{S~}: scope and limitations of the partial-PGF recursion}\label{ssec:A:Stilde}

Setting $\gamma_1=\gamma_2=\theta_1=\theta_2=0$ in the $\widetilde{S}$ balance
equations~\eqref{eq:gen:tilde_balance_interior}--\eqref{eq:gen:tilde_balance_idle}
and forming the PPGF $\widetilde{P}(y,n)=\sum_{n_2=0}^n \widetilde{\pi}(n_2,n)\,y^{n_2}$
yields the inhomogeneous recurrence:
\begin{align}
    \begin{split}
       &\left[ \lambda_1 + \lambda_2 + \mu \right] \widetilde{P}(y,n) \\
       & \quad = \quad(\lambda_1 + \lambda_2 y) \widetilde{P}(y,n-1) + \mu \widetilde{P}(y,n+1)
                  + \mu y^n (1-y) \tpi(n+1,n+1)
    \end{split} \label{eq:A:ppgf_Stilde_dynamics}
\end{align}

The forcing term in~\eqref{eq:A:ppgf_Stilde_dynamics} is controlled by the diagonal
probabilities $\tpi(n+1,n+1)=\pi(0,n+1)$, whose decay we first establish from the exact
$S$-solution.

\begin{lemma}[Geometric decay of the boundary probabilities]\label{lem:A:diag_decay}
    For Model-$A$ the boundary probabilities $\pi(0,n)$ satisfy $\pi(0,n)=O(r^{-n})$ for
    every $r<1/\rho$; equivalently they decay geometrically with rate $\rho$.
\end{lemma}

\begin{proof}
    By definition $P_y(y)=P(0,y)=\sum_{n\ge0}\pi(0,n)\,y^n$ is the generating function of
    $\{\pi(0,n)\}$, and from~\eqref{eq:A:Py} with $\pi(0,0)=\rho(1-\rho)$,
    \begin{equation*}
        P_y(y)=\rho(1-\rho)\,\frac{x^*(y)(1-y)}{x^*(y)-y}.
    \end{equation*}
    Evaluating the kernel quadratic~\eqref{eq:A:xstar_def} at $x=y$ gives
    $\lambda_1 y^2-[\mu+\lambda_1+\lambda_2(1-y)]y+\mu=(\lambda y-\mu)(y-1)$ with
    $\lambda=\lambda_1+\lambda_2$, so $x^*(y)=y$ holds only at $y=1$ and $y=\mu/\lambda=1/\rho$.
    At $y=1$ the simple zeros of $(1-y)$ and $(x^*(y)-y)$ cancel and $P_y$ is analytic with
    $P_y(1)=\rho$; the branch point of $x^*$ lies beyond $1/\rho$. Hence the singularity of
    $P_y$ nearest the origin is the simple pole at $y=1/\rho>1$, so $P_y$ is analytic on the
    disc $|y|<1/\rho$. By the Cauchy coefficient estimates, for every $r<1/\rho$ there is
    $C_r<\infty$ with $\pi(0,n)\le C_r\,r^{-n}$, i.e.\ geometric decay at rate $\rho$.
\end{proof}

\begin{approximation}[Heuristic PPGF on $\widetilde{S}$]\label{apx:A:PPGF_Stilde}
    Neglecting the boundary-diagonal forcing term
    $\mu y^n(1-y)\,\tpi(n+1,n+1)$ in~\eqref{eq:A:ppgf_Stilde_dynamics} --- which is
    exponentially small by Lemma~\ref{lem:A:diag_decay} --- the
    homogeneous recurrence admits the geometric solution
    \begin{equation*}
        \widetilde{P}(y,n) \;\approx\; \rho(1-\rho)\,[\widetilde{y}^*(y)]^n, 
    \end{equation*}
    where
    \begin{equation*}
        \widetilde{y}^*(y) = \frac{(\lambda+\mu) - \sqrt{(\lambda+\mu)^2 - 4\mu(\lambda_1+\lambda_2 y)}}{2\mu}, \qquad \lambda = \lambda_1+\lambda_2. 
    \end{equation*}
    At $y=1$ this recovers $\widetilde{P}(1,n)=(1-\rho)\rho^{n+1}$, consistent with
    the $M/M/1(\lambda,\mu)$ steady-state marginal for $n$ customers waiting.
\end{approximation}

\begin{proof}[Justification of Approximation~\ref{apx:A:PPGF_Stilde}]
The inhomogeneous term $\mu y^n(1-y)\tpi(n+1,n+1)$ prevents the direct application
of a geometric-sequence ansatz. The diagonal probabilities $\tpi(n,n)=\pi(0,n)$
correspond to states where all $n$ waiting customers are class~$2$, and by
Lemma~\ref{lem:A:diag_decay} they decay geometrically in $n$ at rate $\rho$; the
forcing term is therefore exponentially small, and neglecting it yields an
approximation whose accuracy is quantified below. The resulting homogeneous recurrence is
\begin{equation*}
    (\lambda_1 + \lambda_2 + \mu) \widetilde{P}(y,n) = (\lambda_1 + \lambda_2 y) \widetilde{P}(y,n-1) + \mu \widetilde{P}(y,n+1).
\end{equation*}
The characteristic polynomial of this recurrence is
\begin{equation*}
    \mu [\widetilde{y}(y)]^2 - (\lambda_1 + \lambda_2 + \mu) \widetilde{y}(y) + (\lambda_1 + \lambda_2 y) = 0,
\end{equation*}
with roots
\begin{equation*}
    \widetilde{y}^{\pm}(y) = \frac{(\lambda_1 + \lambda_2 + \mu) \pm \sqrt{(\lambda_1 + \lambda_2 + \mu)^2 - 4\mu(\lambda_1 + \lambda_2 y)}}{2\mu}.
\end{equation*}
At $y=1$ the roots are $\widetilde{y}_{+}(1)=1$ and $\widetilde{y}_{-}(1)=\rho$.
The root $\widetilde{y}_{+}=1$ would produce a non-summable geometric series as
$n\to\infty$, so only the negative-sign root is admissible:
\begin{equation*}
    \widetilde{y}^*(y) = \frac{(\lambda_1 + \lambda_2 + \mu) - \sqrt{(\lambda_1 + \lambda_2 + \mu)^2 - 4\mu(\lambda_1 + \lambda_2 y)}}{2\mu}.
\end{equation*}
The geometric ansatz gives $\tP(y,n)=\tP(y,0) \cdot [\widetilde{y}^*(y)]^n$.
Using $\tP(y,0)=\tpi(0,0)=\pi(0,0)=\rho(1-\rho)$ (Corollary~\ref{cor:A:pi0}):
\begin{equation*}
    \tP(y,n)=\rho(1-\rho) \cdot  [\widetilde{y}^*(y)]^n.
\end{equation*}
Since $n$ counts customers \emph{waiting} (the customer in service is not counted),
$n+1$ is the total number in the system, so evaluating at $y=1$:
\begin{equation*}
    \tP(1,n)= \rho(1-\rho) \cdot\rho^n= (1-\rho) \rho^{n+1}, \qedhere
\end{equation*}
which matches the $M/M/1(\lambda,\mu)$ stationary distribution for $n+1$ customers
in the system~\cite{adan2002queueing}.
\end{proof}

\paragraph{Failure mechanism.}
The neglected forcing $\mu y^n(1-y)\,\tpi(n+1,n+1)$ feeds in the diagonal states, in which
all waiting customers are class~$2$. Using $\tpi(n,n)=\pi(0,n)$ and the rate proved
in Lemma~\ref{lem:A:diag_decay}, this term decays in $n$ like $\rho^{\,n}$ (modulated by
$y^n$), whereas the homogeneous solution decays like $[\widetilde{y}^*(y)]^n$. The
approximation is therefore accurate exactly when the forcing decays faster, i.e.\
$\rho<\widetilde{y}^*(y)$ on the relevant range of $y$; this holds when the priority
class carries most of the load ($\rho_1/\rho$ large, $\rho$ moderate), precisely the
empirical validity region of the approximation.

The obstruction is structural and recurs throughout this thesis: an unknown boundary/diagonal
trace --- here the sequence $\pi(0,n)$ --- couples into an otherwise solvable recursion,
exactly as the unknown diagonal $P(z,z)$ enters the Volterra forcing of Models~$B$ and~$X$
(\S\ref{sec:model_b} and Remark~\ref{rem:B:modelX}). A rigorous closure would retain the forcing
and solve~\eqref{eq:A:ppgf_Stilde_dynamics} by variation of constants in $n$, which
reintroduces the very diagonal sequence $\{\pi(0,n)\}$ as an unknown to be determined
self-consistently; see \S\ref{sec:future_work}.
